\documentclass[11pt,a4paper]{article}

\usepackage[T1]{fontenc}
\usepackage[utf8]{inputenc}
\usepackage{lmodern}
\usepackage[scaled=0.92]{inconsolata}
\usepackage[margin=2.6cm]{geometry}
\usepackage{microtype}
\DisableLigatures[-]{encoding = T1, family = tt*}
\usepackage{xcolor}
\usepackage{listings}
\usepackage{booktabs}
\usepackage{tabularx}
\usepackage{enumitem}
\usepackage{caption}
\usepackage[numbers,sort&compress]{natbib}
\usepackage[colorlinks=true,linkcolor=blue!50!black,citecolor=green!35!black,urlcolor=blue!50!black]{hyperref}
\usepackage[capitalise,noabbrev]{cleveref}

\definecolor{kw}{RGB}{0,60,140}
\definecolor{doc}{RGB}{110,110,110}
\definecolor{str}{RGB}{150,40,40}
\definecolor{bg}{RGB}{248,248,246}

\lstdefinelanguage{ulang}{
  morekeywords={function,procedure,data,type,export,implementation,builtin,import,
    match,case,end,if,then,else,let,in,out,ref,local,assign,call,raise,raises,try,catch,
    skip,while,do,for,yield,lambda,test,assert,lemma,requires,ensures,pattern,infix,
    prefix,postfix,and,or,not,example,proof,induction,using,have,by,decreases,uses,correct,old,with},
  sensitive=true,
  morecomment=[l]{//},
  morestring=[b]",
  morestring=[b]',
}
\lstset{
  language=ulang,
  basicstyle=\ttfamily\small,
  keywordstyle=\color{kw}\bfseries,
  commentstyle=\color{doc},
  stringstyle=\color{str},
  columns=fullflexible,
  keepspaces=true,
  showstringspaces=false,
  backgroundcolor=\color{bg},
  frame=none,
  xleftmargin=0.8em,
  framexleftmargin=0.4em,
  aboveskip=0.8em,
  belowskip=0.6em,
  upquote=true,
}
\lstnewenvironment{diag}{\lstset{language={},basicstyle=\ttfamily\footnotesize,keywordstyle=,stringstyle=}}{}
\newcommand{\code}[1]{\texttt{\detokenize{#1}}}
\newcommand{\file}[1]{\texttt{#1}}
\newcommand{\planned}{\textsuperscript{\textsf{\scriptsize P}}}

\setlist{itemsep=0.2em,topsep=0.4em}
\captionsetup{font=small,labelfont=bf}

\title{\textbf{ulang: An Agent-First Programming Language}\\[0.4em]
\large Design Rationale, Current State, and the Road Ahead}
\author{Gidon Ernst}
\date{November 2026 \quad---\quad internal design paper, not versioned}

\begin{document}
\maketitle

\begin{abstract}
\noindent
Writing code has become nearly free; reading, checking, and changing it safely has not.
When the main author of a program is an AI agent, the agent's typing does not cost much.
What costs is two other things: the agent learning the language from the few pages in its
context, and someone (a person, a tool, or another agent) verifying what it wrote.
This paper describes \emph{ulang}, a small language designed under that assumption, and
implemented in itself: a parser, a type checker with Hindley--Milner inference and checked
exceptions, an interpreter, a compiler to C, and a prover for lemmas about pure functions,
about 19\,000 lines of ulang that compile themselves to a fixed point. We state the design
principles that follow from the agent-first assumption: locality of every rule, one way to
write each construct, a canonical layout enforced by the checker, effects visible in
signatures and at calls, atomic effects under exceptions, and failure that is never silent.
We then describe the language as it will stand once the plans now in progress are in. These
are diagnostics as a machine-readable protocol with hints for the habits of other
languages, examples that are part of signatures and run as tests, lemmas that are tested
on enumerated values before they are proven, a structured proof language over an e-graph
simplifier, lint rules that are equalities the prover establishes, and separate compilation.
For each property that makes code cheap to verify, we say how the design provides it and
where it falls short. We close with the tensions that remain and the experiments that would
settle them.
\end{abstract}

\tableofcontents
\bigskip

%=====================================================================
\section{Introduction}
%=====================================================================

Most conventions of programming languages grew around human authors. Short notation saves
typing, implicit conversions save thought, and the many ways to write a loop let each author
keep a personal style. Readers pay for all of this, in context they must recover and in
alternatives they must recognise. While writing was the expensive part, the trade was
reasonable.

Coding agents change both sides of that trade. An agent writes thousands of lines without
fatigue, so a keyword more or less costs it nothing. On the other side, agents make a kind
of mistake that conventional languages hide well: well-formed, plausible code that misreads
the intent. A shadowed name, an implicit fallback, a forgotten case, an effect that happens
somewhere it should not. Whoever checks the code (a person, a reviewing model, a tool) now
dominates the cost. Following the companion sketch on non-functional properties of AI code,
we take \emph{cost to verify} as the measure of structural quality: how much a verifier must
know and check to be confident that a unit is correct and that a change to it is
safe.

A new language meets a second problem, one particular to agents: there is no ulang in any
model's training data. An agent that writes ulang learns it inside one session, from a page
of documentation and from the messages of the checker. Two things follow.
\begin{itemize}
\item \textbf{The cost of a feature is the length of the paragraph that explains it to an agent.}
  A feature that needs a long explanation is expensive however much it saves in typing.
\item \textbf{The checker is the teacher.} Every error message is a lesson. An error that
  names only the token expected (\code{end expected but x found}) teaches little. An error
  that names the construct to write instead teaches the language.
\end{itemize}

ulang is an experiment in taking these premises seriously. It has pure functions and
imperative procedures, algebraic datatypes, Hindley--Milner inference~\cite{milner1978}
within declarations, checked exceptions, unbounded integers, ordered containers,
comprehensions, modules with explicit exports, and lemmas with an automatic prover. It is
implemented in itself, and the implementation has been developed largely with coding agents,
so the language is also its own main test of whether it suits them. A file of trip-ups,
\file{unclear.md}, records each place where an agent guessed wrong, so that the documentation
can be fixed from it.

\paragraph{Scope of this paper.}
Several plans are in progress as notes in the repository: diagnostics (\file{diagnostics.md}),
examples and tested lemmas (\file{examples.md}), lint (\file{lint.md}), the prover and its
proof language (\file{prove.md}, \file{prover.md}), separate compilation
(\file{archives.md}), and a native backend (\file{assembly.md}). We describe the language
\emph{as if those plans were done}, since that is the design to discuss. Where something is
not yet implemented we mark it with a superscript~\planned, and \cref{sec:status} gives the
status of each part. Items of \file{TODO.md} that no plan covers yet are discussed as open
questions in \cref{sec:discussion}.

\paragraph{Contributions.}
\begin{enumerate}
\item A set of design principles for a language whose authors are agents, derived from cost to
  verify and in-context learnability (\cref{sec:principles}).
\item A language that applies them throughout, and a self-hosting implementation as evidence
  that the result is practical (\cref{sec:language,sec:impl}).
\item A view of diagnostics as a protocol between checker and agent rather than text for
  people (\cref{sec:feedback}).
\item A \emph{ladder of evidence} that runs from types through examples, enumerated tests,
  and the default proof up to structured proofs and contracts, where each rung reuses the one
  below (\cref{sec:spec}).
\item Lint as extraction from an e-graph: a suggested rewrite is a cheaper representative of
  a class that the prover's rules populate, so every suggestion is a proven equality
  (\cref{sec:lint}).
\item An account of what remains in tension (\cref{sec:discussion}).
\end{enumerate}

%=====================================================================
\section{Design Principles}\label{sec:principles}
%=====================================================================

The principles below are those of \file{DESIGN.md}, sharpened by the agent-first premise.
Each one names the property of cost to verify that it serves; the properties are those of
the sketch on non-functional properties (locality, explicitness, non-redundancy,
proportionality, fail-loud behaviour, grounding, traceability, uniformity, and proportionate
defensiveness).

\begin{description}[style=nextline,leftmargin=1.2em]
\item[P1. Locality.]
  Every rule can be decided from the code at the point of use, its arguments, and the
  signatures it refers to. It never depends on other uses in the same command, on line
  breaks, or on preferences among candidates. Overloading resolves uniquely or fails
  (\cref{sec:overloading}). Imports are not transitive. User-defined notations are the one
  deliberate exception. \emph{Serves: locality, explicitness.}

\item[P2. One way to write each thing.]
  Two syntaxes for one construct are two things to learn and a chance to mix them up. Lists
  are \code{List(1, 2)}, never \code{[1, 2]}. A chain of conditions is a nested \code{if},
  never \code{elif}. Rejected alternatives are recorded with their reasons, so that no later
  session adds them back. \emph{Serves: uniformity.}

\item[P3. Keystrokes are cheap; tokens read are not.]
  A convenience that only saves typing is not worth a rule. Every statement starts with its
  keyword (\code{local}, \code{assign}, \code{call}, \code{raise}, \code{skip}), and every
  argument of a call carries its mode (\code{in}, \code{out}, \code{ref}). This verbosity is
  low in entropy: a reader skims the keywords and gains certainty about the shape. Semantic
  redundancy, such as restating in a comment what a name says, is what costs context, and the
  style rules forbid it. \emph{Serves: explicitness.}

\item[P4. The layout is canonical, and the checker enforces it.]
  \code{ulang format} writes the only accepted layout, and \code{check} rejects any other.
  An agent cannot introduce style drift, and a diff never reformats code it does not change.
  The formatter is also a test: every file that checks has passed a round trip through
  parser and printer. \emph{Serves: proportionality, uniformity.}

\item[P5. Comments are syntax.]
  Prose is a \emph{doc} attached to a module, a command, a constructor, or a \code{case}.
  It is kept in the tree and printed back, so no tool loses or moves it. A comment at the end
  of a line or in the middle of a body is a parse error: the code there needs a name instead.
  Docs that make claims become \emph{examples}, which are checked (\cref{sec:examples}).
  \emph{Serves: traceability, grounding.}

\item[P6. Effects show in signatures and at calls.]
  The kind of a declaration is a one-bit effect system: functions are pure expressions,
  procedures are statements. Data flow shows in the modes of formals and arguments. What a
  procedure may raise is part of its type. \emph{Serves: locality, explicitness.}

\item[P7. Effects are atomic.]
  A procedure works on copies of its outs and refs, which reach the caller only when it
  returns. The handlers of a \code{try} see the variables as they were before it. There is
  no state that is half updated when an exception escapes. \emph{Serves: locality, fail-loud.}

\item[P8. Failure is never silent.]
  Functions cannot \code{raise}; they fail only by the exceptions of builtins, which are
  defects. A failed \code{assert} cannot be caught. Matches are exhaustive. Lookups give
  options, and a default is named at the call (\code{value_or}). Generators do not skip
  elements whose shape does not match. A handler of an exception that its body cannot raise is
  an error. \emph{Serves: fail-loud, proportionate defensiveness.}

\item[P9. Determinism.]
  Integers are unbounded. Maps and sets are ordered by a total order on all values, so
  iteration never depends on hashing. The interpreter and compiled code agree, with the texts
  of shared runtime errors in one place. The prover works within named budgets and reports
  when one runs out. Tests of properties enumerate small values rather than drawing random
  ones. \emph{Serves: grounding, reproducibility of every result an agent sees.}

\item[P10. A small type system, with contracts for the rest.]
  Types are the decidable, always-checked part of a specification: inference inside bodies,
  declared types at boundaries, no higher kinds, no computation on types. Everything more
  expressive goes into examples, lemmas, and contracts, which are checked by tests or the
  prover (\cref{sec:spec}).

\item[P11. Every departure from a common habit has a hint.]
  This principle is new to this paper. Agents arrive with the habits of Python, Rust, and
  TypeScript, and ulang rejects some of them on purpose. A rejection is a learning cost,
  which a hint that names the replacement pays back
  (\cref{sec:hints}). A departure that cannot be given such a hint is a sign that it is
  arbitrary.
\end{description}

\noindent
Two consequences shape the rest of the design. First, \emph{the parser desugars nothing}.
Comprehensions, \code{for} loops, patterns in lambdas, field updates, and pattern aliases
stay as written until the checker translates them, so the formatter prints what the source
wrote and an error is reported against what the agent wrote. Second, \emph{tools are part of
the language}. Formatting, resolution of names, refactorings, lint, and the prover read and
write the same syntax tree, and the commands of one binary are the interface through which an
agent works.

%=====================================================================
\section{The Language}\label{sec:language}
%=====================================================================

\subsection{A first example}

\Cref{fig:bank} shows a complete program. It checks, its tests pass, the lemma is proved
automatically, and its layout is the one \code{format} writes.

\begin{figure}[t]
\begin{lstlisting}
data Account :=
  case Account(owner: Str, balance: Int)

data Overdrawn :=
  case Overdrawn(missing: Int)

// takes amount from acct, which stays as it was if the balance does not cover it
procedure withdraw(in amount: Int, ref acct: Account) raises Overdrawn :=
  if acct.balance < amount then
    raise Overdrawn(amount - acct.balance)
  end
  assign acct := acct with balance := acct.balance - amount

procedure transfer(
  in  amount: Int,
  ref from: Account,
      to: Account
) raises Overdrawn :=
  call   withdraw(in amount, ref from)
  assign to := to with balance := to.balance + amount

function totals(accts: List[Account]): Map[Str, Int] :=
  Map(for a in accts yield a.owner => a.balance)

lemma size_append[A](xs: List[A], ys: List[A]) :=
  ensures size(xs ++ ys) == size(xs) + size(ys)

test kept() :=
  local a := Account("ada", 10), b := Account("bob", 0)

  try
    call transfer(in 15, ref a, b)
  catch
    case Overdrawn(_) =>
      skip
  end
  assert a.balance == 10 and b.balance == 0
\end{lstlisting}
\caption{A program in ulang. \code{ulang test} prints \texttt{test kept passed} and
\texttt{lemma size\_append proved}.}
\label{fig:bank}
\end{figure}

The example shows most of the core. Datatypes declare every constructor after \code{case},
and every field has a name. Fields are read as \code{e.f} and updated as
\code{e with f := v}. Procedures take formals grouped by mode, and calls repeat the modes.
Exceptions are values of datatypes, and a procedure says what it raises. The handler in
the test sees \code{a} and \code{b} as they were before the \code{try}, so the final
assertion holds even though \code{withdraw} was reached. A comprehension builds a map, and a
lemma states a property of a library function. The doc above \code{withdraw} is part of the
syntax.

\subsection{Functions, procedures, and modes}

Functions are expressions: \code{let p := a, q := b in e} binds patterns in parallel,
\code{if c then a else b} always has its \code{else}, and \code{match} is exhaustive and
checked for unreachable cases. Procedures are statements:
\code{local}, \code{assign} (simultaneous: \code{assign a := b, b := a} swaps),
\code{call}, \code{if}/\code{while}/\code{for} with \code{end}, \code{try}/\code{catch},
\code{raise}, \code{assert}, and \code{skip}. There is no \code{return} and no
\code{break}. A procedure gives its results by its outs, and a loop ends by its condition or
an exception. Variables are read only where they are assigned on every path, and outs must be
assigned on every path that does not raise.

Both kinds are values. \code{lambda x, (a, b) => e} is a function whose formals are
irrefutable patterns. \code{lambda in x, out y => ... end} is a procedure of type
\code{(in A, out B) raises E}, which sees the variables around it by value and assigns only
its own outs and refs. Procedure values abstract over effectful code: the parser of ulang
writes its separated lists once, as a procedure that takes the parser of an element.

No formal, local, or pattern variable may hide a local around it (\code{HiddenLocal}).
Shadowing is a classic source of plausible-looking bugs, and an agent rarely means it.

\subsection{Exceptions, checked and inferred}\label{sec:raises}

A raised value is of a datatype without type parameters. \code{raise} is a statement, so
only procedures raise. What a procedure may raise is checked after its package:
\begin{itemize}
\item A statement raises the types of its \code{raise}, the \code{raises} of the procedures it
  calls (also through variables of procedure types), and what a \code{try} lets through: the
  types that its unguarded handlers do not cover, decided by the exhaustiveness check.
\item Exported procedures and those with \code{raises} raise nothing they omit. Private
  procedures without \code{raises} infer theirs, as the least fixpoint over all of them.
  Annotation is required only at boundaries, where a reader needs it.
\item Procedure types carry their raised types as a set, compared exactly, without
  subtyping. A lambda gets an unknown that its place of use fixes.
\item A handler of a type that its body cannot raise is an error. So is a handler that
  nothing in its \code{try} reaches. Dead defensive code is not allowed to accumulate.
\end{itemize}
The exceptions of builtins (\code{DivisionByZero}, \code{OutOfBounds}, \code{NoCharacter},
\code{MissingField}) are not counted, since functions cannot declare them. They are defects of
the program, and the planned \code{requires} of functions rules them out
statically~(\cref{sec:contracts}).

\subsection{Datatypes, fields, and pattern aliases}

Datatypes are sums of records. A field name may repeat across constructors with the same
type, and an accessor applied to a constructor without the field raises \code{MissingField}.
\emph{Pattern aliases} name a shape at the level where code works with it:
\begin{lstlisting}
pattern ECall(name: Name, typ: Type, args: List[At[Expr]]): Expr :=
  EApp(AOk(_, EFun(name, typ)), args)
\end{lstlisting}
An alias is typed once in its module and expanded at its uses, so exhaustiveness and the
backends see only constructors. An alias without \code{_} also builds values. Unlike view
patterns, aliases call no function, so exhaustiveness stays decidable and matching always
terminates.

\subsection{Containers as forms}

The containers are written as \emph{forms} named like their types: \code{List(a, b; rest)},
\code{Set(a, b)}, \code{Map(k => v)}, \code{Array(a, b)}. Comprehensions use the same form
with generators and guards from left to right:
\begin{lstlisting}
List(for x in xs for y in ys if x < y yield (x, y))
Map(for k => v in m if v > 0)
\end{lstlisting}
The form names what is built, so the result never depends on the expected type. That is the
problem of Scala's \code{CanBuildFrom}, and it would violate locality. The checker
translates forms into the folds of the container's module, chosen by the container's type.
Generators range over lists, sets, maps, arrays, buffers, and options alike, and
\code{for x in c if q do ... end} is the matching statement. Their patterns are irrefutable:
a refutable pattern would either drop elements silently or fail at run time, and both were
rejected. The shared names (\code{size}, \code{contains}, \code{get}, \code{remove},
\code{foldl}, \code{foldr}, \code{++}, \code{--}) take their arguments in the same order for
every container, so knowing one container means knowing them all.

\subsection{Overloading without preferences}\label{sec:overloading}

Within a module, a name overloads only by arity. Across modules it overloads by type.
A use infers its arguments, keeps the candidates of its arity whose instance unifies with them
and with the expected type, and commits if exactly one remains. Pending uses are solved to a
fixpoint at the end of the command. Since the substitution only grows, a candidate that does
not fit never fits again, so the result does not depend on the order of solving. There are no
preferences, no defaulting, and no implicit conversions. An ambiguous use is an error that
lists the candidates with their types:
\begin{diag}
e3.u: type error: empty is ambiguous: prelude/set.empty: () -> Set[A],
  prelude/map.empty: () -> Map[K, V], prelude/buffer.empty: () -> Buffer[A] at 2.14
\end{diag}
The planned hint turns this into the annotation that decides it,
\code{local s: Set[Int] := empty()}\planned.

The \emph{hooks} \code{compare}, \code{==}, and \code{show} follow the same rule. A function
of the right shape in the module of a datatype becomes that type's order, equality, or text,
and the builtins reach it through the type descriptor of each value, also inside other values
and in generic code. Only abstract types may declare \code{==}. Comparing values of an
abstract type without an order raises \code{Unordered}, since the shape of an AVL tree is no
order.

\subsection{Modules and packages}

A file with \code{export} and \code{implementation} sections is a module, and the modules of
a directory form a package. Members see each other's exports, also cyclically. Packages are
imported by their path from the root of the project, \code{import "ulang/syntax"}, and
imports are not transitive, so the import lines of a file name every package its names come
from. Export sections hold signatures without bodies, which the implementation repeats
exactly, so an agent that reads an export section sees the whole interface. Abstract types
hide their constructors. Relative imports, re-exports, and subpackages private to a parent
were each considered and rejected: each would add a second way to reach the same declaration,
or a second spelling of the same package.

\subsection{Tests as part of the language}

Tests are commands, \code{test name() := ...}, next to the code they test, where they see
private declarations. \code{assert raises p do ... end} expects an exception that matches
\code{p}. Errors that a phase must report are themselves tests: a case is a source text in a
test of \file{ulang/test/errors/}, run through the phases, with the expected exception
matched by \code{assert raises}, including its position. Every file of the repository passes.
There are no ``expected to fail'' files, so ``all files pass'' is the one rule to check.

%=====================================================================
\section{Feedback as a Protocol}\label{sec:feedback}
%=====================================================================

An agent works in a loop: write, check, read the errors, fix. The quality of that loop
depends on three things: how many errors one run reports, how precisely each is placed,
and whether it says what to do.

\subsection{Positions on every part}

Every part of an expression, statement, and pattern is a value of
\begin{lstlisting}
data At[A] :=
  case AOk(pos: Pos, that: A)
  case AErr(pos: Pos, that: A, error: Diagnostic)
  case AUnparsed(pos: Pos, text: Str, error: Diagnostic)
\end{lstlisting}
Positions are never stripped. They reach the checked program, the runtime errors\planned, and
the open goals of the prover\planned. Code that does not care reads \code{x.that}. A type error is
reported at the innermost part around it, and a part of the wrong type is a \code{Mismatch}
at that part, with the type expected, the type found, and the role of the part:
\begin{diag}
e2.u: type error: expected Int but found Str in argument 1 of the call at 2.3
e2.u: type error: match is not exhaustive, missing None at 5.3
\end{diag}

\subsection{All errors in one run}

The checker records each error and continues with a fallback: a fresh type for an
expression, a statement counted as assigning what it assigns, a pattern as a wildcard. So one
run reports the errors of all commands of the first package that has some, at most 20 per
file, without cascades from the fallbacks. The parser resumes at the next command, and in the
planned recovery also at the next statement and argument\planned. What does not parse becomes an
\code{AUnparsed} part, and the commands that parse are checked.

Errors are data from start to end. Each is a \code{Diagnostic(kind, code, module, pos,
message, details)}, where \code{code} is the constructor of the error and \code{details}
holds pairs such as \code{expected}, \code{found}, \code{candidates}, and \code{missing}. The
driver prints them as one JSON object per line\planned, the only output format, and the exit
status separates source errors, usage errors, failed tests, and internal errors\planned. The
editor extension, scripts, and agents all read the same records.

\subsection{Hints for the habits of other languages}\label{sec:hints}

The checker and the parser often know the fix. The most frequent mistakes of agents writing
ulang are habits from other languages, and today they get only the token that was expected:
\begin{diag}
e1.u: parse error: declaration expected but x found at 5.3     -- from:  x := 1
\end{diag}
The plan attaches a hint to each such habit\planned\ (\cref{tab:hints}). Each hint has a test
case. The \code{end} of an \code{if} after \code{else} is reported at the \code{if} it closes,
found by indentation, instead of at the end of the input. Suggested edits, such as the
missing cases of a \code{match} or the annotation that decides an ambiguous use, are listed
with the diagnostic and applied by \code{--apply} as refactorings
(\cref{sec:refactor}).

\begin{table}[t]
\centering\small
\begin{tabularx}{\textwidth}{@{}lX@{}}
\toprule
\textbf{Habit} & \textbf{Hint names} \\
\midrule
\code{x := e} as a statement & \code{assign x := e} \\
\code{f(x)} as a statement & \code{call f(in x)} \\
\code{return e} & a function is its expression; a procedure assigns its outs \\
\code{[1, 2]}, \code{x :: xs} & \code{List(1, 2)}, \code{List(x; xs)} \\
\code{elif} & \code{else if c then ... end}, with its own \code{end} \\
\code{y += 1} & \code{assign y := y + 1} \\
\code{s + "x"} on strings & \code{++} \\
procedure used as a function & \code{call f(in ..., out r)} \\
\code{local} after a statement & locals are declared at the start of a body \\
unknown name & the nearest visible names, and packages not imported that export it \\
argument order & a permutation of the arguments that fits a candidate \\
\bottomrule
\end{tabularx}
\caption{Planned hints for habits that agents bring from other languages. The current
messages for the first two are \texttt{declaration expected}; for the string case, a
cascade of three mismatches.}
\label{tab:hints}
\end{table}

\subsection{Seeing what the checker resolved}

Resolution by type is invisible in the text: \code{size} in a file that imports \code{set}
and \code{map} may be any of three. The checked program keeps a position for every part, so a
use can be looked up by its position. The tool offers \code{resolve <file> <L.C>} (qualified
name, type, and declaration)\planned, \code{refs} (all uses)\planned, and \code{sig <module>}
(the export section with its examples)\planned, each also as hover in the editor. These cost
almost nothing to build, and they turn ``read the imports and guess'' into one lookup.

\subsection{Refactorings that are safe under overloading}\label{sec:refactor}

Refactorings are changes of the parsed tree, printed by the formatter, which keeps docs since
they are syntax\planned. A refactoring is refused unless the file is formatted, the result is checked
again, and it is refused if any other use resolves differently than before. That last check
makes renaming safe under overloading without any reasoning about overloading. The planned
order is rename, filling in the missing cases of a match, adding a field with its
constructions and patterns, and changing a signature with its calls. Extracting and inlining
are left out, since textual edits do those well enough. Typed edits are also the strongest
lever for proportionality: a change described as ``add a field'' touches exactly what that
change needs.

%=====================================================================
\section{A Ladder of Evidence}\label{sec:spec}
%=====================================================================

The common mistakes of agents are logic errors, not type errors. ulang therefore puts
specification in the language, as a sequence of rungs. Each rung costs a little more to write
than the one below, gives stronger evidence, and reuses the machinery below it
(\cref{tab:ladder}).

\begin{table}[t]
\centering\small
\begin{tabularx}{\textwidth}{@{}llX@{}}
\toprule
\textbf{Rung} & \textbf{Written as} & \textbf{Checked by} \\
\midrule
types & signatures & the checker, always \\
examples\planned & \code{example gcd(12, 18) == 6} & run as tests, \code{gcd.example.1} \\
tests & \code{test t() := assert ...} & interpreter and compiled code \\
enumerated lemmas\planned & \code{lemma l(xs: List[A]) := ensures ...} & all small values, \code{l.lemma} \\
default proof & the same lemma & simplification, else induction \\
structured proof\planned & \code{proof induction xs case ... end} & the simplifier, step by step \\
contracts\planned & \code{requires}, \code{ensures}, \code{correct p} & proof obligations of calls and loops \\
\bottomrule
\end{tabularx}
\caption{The ladder of evidence. A status per obligation (proven, tested, or open) and
\code{ulang audit}, which lists what depends on an assumption, make the rung of every claim
visible\planned.}
\label{tab:ladder}
\end{table}

\subsection{Examples in signatures}\label{sec:examples}

A signature gives types. A doc gives prose that nothing checks and that goes stale. Tests sit
elsewhere, private to their module. An agent that reads an export section therefore never sees
a concrete call with its result. Examples close this gap\planned:
\begin{lstlisting}
export

// the greatest common divisor, nonnegative
function gcd(a: Int, b: Int): Int
  example gcd(12, 18) == 6
  example gcd(0, -5) == 5

procedure divmod(in a: Int, b: Int, out q: Int, r: Int)
  example divmod(in 7, 2, out 3, 1)
  example divmod(in 7, 0, out _, _) raises DivisionByZero
\end{lstlisting}
An example is a closed expression that must apply its own declaration, as overloading resolves
it. So \code{example size(List(1)) == 1} on \code{set.size} is an error, \code{ForeignExample}.
Examples in an export section see only what an importer sees, so the example of an abstract
type builds its values through the exported functions. The checker translates examples into
tests named \code{gcd.example.N}, so runner, interpreter, compiler, and runtime need no
change. Examples answer the most common question about a signature whose types leave it open:
the order of arguments of \code{slice}, \code{split}, or \code{foldr}. They are the
documentation that cannot go stale, and the first values that generated tests try.

\subsection{Lemmas, tested before they are proven}

A lemma states a property for all formals that satisfy its assumptions:
\begin{lstlisting}
lemma take_drop[A](n: Int, xs: List[A]) :=
  ensures take(n, xs) ++ drop(n, xs) == xs
\end{lstlisting}
Lemmas are checked as \code{Bool} expressions over their formals, so they carry resolved names
like any code. Before a lemma is proven, the checker generates a test \code{l.lemma} that
enumerates all values of the formals up to a depth, in the style of
SmallCheck~\cite{runciman2008smallcheck}\planned. It is deterministic without a seed, the
first counterexample is among the smallest so no shrinking is needed, and small cases are
where lemmas about lists and integers fail. Type parameters are instantiated to \code{Int}.
A lemma whose assumptions no small case meets fails as \code{Vacuous}. Enumerators for each
datatype are generated as functions that recur on the depth, so they terminate, and they take
the enumerators of type parameters as functions. Abstract types of users are not enumerated,
since their constructors do not keep their invariants.

The order matters for agents. A false lemma is refuted in milliseconds with a concrete
counterexample, before anyone spends effort on a proof. That gives a property the same tight
feedback loop as a type error.

\subsection{The default proof and structured proofs}

\code{ulang prove} tries each lemma by simplification, and failing that by induction on one
of the formals, with the lemmas proven before as facts. When neither closes the goal it
reports the simplified goal, which serves as the error message of a proof. The planned proof
language\planned\ adds as little as possible, with the syntax habits of the rest of ulang:
every step starts with its keyword, blocks end with \code{end}, and cases are written as in
\code{match}.
\begin{lstlisting}
lemma reverse_onto_append[A](xs: List[A], ys: List[A]) :=
  ensures reverse_onto(xs, ys) == reverse(xs) ++ ys
proof
  induction xs
    case List(x; rest) =>
      using reverse_onto_append(rest, List(x; ys)), reverse_onto_append(rest, List(x))
  end
\end{lstlisting}
The steps are \code{using} (facts, lemmas, instances, or definitions), \code{have h := e}
with an optional \code{by}, \code{induction x}, \code{match}, \code{if}, and closing by an
instance. Simplification runs before each step and at the end of each sequence. Cases that it
closes may be omitted, and a written case that it would close is allowed, so proofs survive
a stronger simplifier. The induction hypothesis is the lemma itself, by its name, at any
arguments, so generalisation comes for free. Open goals are reported at the step that leaves
them, so they appear in the editor like type errors.

\subsection{Contracts}\label{sec:contracts}

Functions get \code{requires}, which states their domain. That includes the builtins:
\code{requires d != 0} of \code{/} turns \code{DivisionByZero} from a defect at run time into
an obligation at the call. Facts about the results of functions are lemmas. Procedures get
\code{requires} and \code{ensures}, where \code{old(x)} is the initial value of a ref, and a
lemma \code{correct p(in n, out s)} claims that the body meets the contract\planned:
\begin{lstlisting}
procedure sum(in n: Int, out s: Int) :=
  requires n >= 0
  ensures  2 * s == n * (n - 1)
  ...

lemma sum_correct(n: Int): (s: Int) :=
  correct sum(in n, out s)
proof
  at while i < n
    invariant 0 <= i and i <= n and 2 * s == i * (i - 1)
\end{lstlisting}
Calls use contracts, not bodies. Each body is proven assuming all contracts, including its
own, which is sound for partial correctness by induction on the depth of calls
(Hoare's rule for recursive procedures~\cite{hoare1971procedures}). Proofs therefore have no
order and no cycles. Loops are specified in the proof, not in the code, either by an
invariant or by a \emph{loop contract} \code{requires P ensures Q} that treats the loop as a
recursive procedure~\cite{ernst2022loops,tuerk2010loops}. The proof refers to points of the
code by quoting them (\code{at while i < n}), matched modulo positions, which reads without
tools and survives most edits. Keeping invariants out of the code keeps the code readable for
the agent that maintains it. The proof then becomes a separate artefact, as tests already are.

Two properties make contracts suited to agents. First, every contract is an obligation even
without a proof: the verifier tries a default proof and reports where it stops, so the
workflow is to write the contract, see where it gets stuck, and add a lemma. Second, an
unproven contract makes the proofs that rely on it conditional, and the tool says so
(\texttt{sum proved, assuming q}). Evidence is never overstated.

%=====================================================================
\section{The Prover}\label{sec:prover}
%=====================================================================

There are no external solvers. The prover is about 2\,500 lines of ulang in
\file{ulang/prove/}, working on checked programs, so its terms carry the names the checker
resolved. Its builtins are pattern aliases over those names.

\paragraph{Sequents.}
Goals, facts, lemmas, and induction hypotheses are nested sequents
$\forall \bar x.\; A_1 \wedge \dots \wedge A_k \Rightarrow
  (\exists \bar y_1.\, C_{11} \wedge \dots) \vee \dots$
whose parts are again sequents or atoms. A lemma \emph{is} a sequent: proving it opens it,
and using it adds it unopened as an assumption. Normalisation makes only steps that lose
nothing (one-point rules, existentials of assumptions to fixed variables), so it never turns
a provable goal into an unprovable one.

\paragraph{Terms.}
Binders are de Bruijn indices and the variables of the context are levels, so a term of an
outer sequent is valid unchanged in inner ones. The state stored at a sequent is a persistent
value shared by those nested in it, and splits need no undo.

\paragraph{Contextual simplification.}
Each part of a sequent is replaced by one that is equivalent in its context. Assumptions
become state of the decision procedures (ground atoms), propagating clauses, or rewrite
rules (oriented equations). Recursive functions unfold where their outermost \code{match}
or \code{if} is decided. Linear arithmetic is canonical sums with Fourier--Motzkin
elimination and integer tightening. Constructors are injective, distinct, and acyclic.
Splits on \code{if} and \code{match} are kept only if every branch closes, so the printed goal
is always a single goal, equivalent to the original.

\paragraph{The e-graph core\planned.}
The current core is a substitution with ground completion. It works, but it is fragile in
specific ways. Fuel numbers are spread over the code, matching is syntactic, the first
matching rule wins, and an equation that the term order cannot orient, such as commutativity,
closes only its exact instance. The plan replaces the term layer with a persistent e-graph
with bounded equality saturation~\cite{willsey2021egg}, in the style of
CC(X)~\cite{conchon2008ccx,nelson1979,shostak1984}. Constructors and literals become class
analyses, linear arithmetic is canonised over classes with entailed equalities fed back, and
facts, lemmas, and the equations of definitions become rules applied by e-matching, all at
once per round. Unfolding becomes guarded e-matching: \code{size(Cons(x, xs))} fires
whenever the argument's class contains a \code{Cons} node, whatever produced that knowledge.
The term order stays, but only to choose the representative that is \emph{printed}, so
provability no longer depends on orientation while the open goal stays predictable.

\paragraph{Budgets, not fuel.}
All limits move into one record of named numbers: nodes, rounds, nested conditions, splits.
An open goal that ran out of budget says so, e.g.\ \texttt{(budget: nodes)}, so an agent
knows whether to add a fact or to raise a limit. Determinism (P9) holds throughout: the same
input gives the same goal, with no timeouts.

\paragraph{Decision procedures, by value for ulang's lemmas.}
The planned order is: evaluating the order on values (\code{compare} of constructor terms);
\code{compare} as a theory; ground clauses with bounded splitting; \code{Bool} and other
finite types as splits; constructor exhaustiveness (\code{xs != List()} gives
\code{xs == List(y; ys)}); division and remainder by literals; integer completeness by the
Omega test~\cite{pugh1991omega}; and ring normalisation. One hazard has to come first:
\code{Set} and \code{Map} are AVL trees whose shapes differ for equal values, so their
constructors must stay uninterpreted for \code{==} before \code{compare} is evaluated.
Otherwise the prover would refute true equations about sets.

\paragraph{Termination.}
Unfolding a function that does not terminate can prove false: \code{f(x) = f(x) + 1}. The
checker will accept recursion in functions only if it is structural or decreases a declared
measure, and will advise turning the function into a procedure otherwise\planned. Until then lemmas
stay private, and the usage message states the risk.

%=====================================================================
\section{Lint as Proven Rewrites}\label{sec:lint}
%=====================================================================

The design already provides much of what makes code cheap to verify. \Cref{tab:props} lists
each property, what the design gives, and what a check still has to find.

\begin{table}[t]
\centering\small
\begin{tabularx}{\textwidth}{@{}l>{\raggedright\arraybackslash}X>{\raggedright\arraybackslash}X@{}}
\toprule
\textbf{Property} & \textbf{By construction} & \textbf{By a check (planned)} \\
\midrule
locality & pure functions, no globals, non-transitive imports, atomic effects & unused imports as errors; \code{uses os}; unused declarations \\
explicitness & nothing desugared in the parser; \code{call}, \code{assign}, modes at calls & \code{resolve}; agreement of declared hooks \\
non-redundancy & one library, shared names across containers & duplicate definitions; library code spelled out \\
proportionality & canonical format; typed refactorings & a diff of declarations modulo equality \\
fail-loud & no \code{raise} in functions, uncatchable \code{assert}, exhaustive matches, options & swallowed handlers; masked domain errors \\
grounding & every name resolves or errs; types fix most argument orders & hints for near misses; examples \\
traceability & --- & unused locals, constructors; coverage by tests \\
uniformity & one way of each construct & rewrites to the canonical form \\
defensiveness & unreachable cases, unreachable handlers & conditions decided by context; redundant guards \\
\bottomrule
\end{tabularx}
\caption{The properties of cost to verify, and how ulang meets them.}
\label{tab:props}
\end{table}

The interesting part is how the rewrites are found. A rewrite is an equality plus a choice of
the simpler side, and the e-graph of the prover separates the two\planned:
\begin{itemize}
\item \textbf{Saturation finds the equalities}, over the definitions of the library, the
  proven lemmas, and the facts of the path (conditions of \code{if}s, shapes of
  \code{case}s).
\item \textbf{Extraction chooses} the cheapest representative, where the cost for lint is the
  number of tokens of source.
\end{itemize}
Each expression of a body enters the e-graph by the prover's translation, and lint keeps a
side table from classes to the parts of the source in them. A lint is reported where a class
has a printable representative cheaper than the part by a margin. A representative is
printable only if it is a source part or the right side of a rule as written. So a hand-written
\code{match} on an option becomes \code{value_or(o, d)} because the definition of
\code{value_or} put it in the same class, and \code{List(for x in xs)} becomes \code{xs}
because a lemma of the library says so. Whoever states a lemma in the library has thereby
written a lint rule, which a change of the definition breaks by failing the proof.

Three safeguards keep suggestions honest. A replacement is checked again, and dropped if one
of its names resolves differently in the module. A representative whose rule would evaluate
an argument that the original reached only under a branch is printable only if that argument
is a plain term, which preserves strictness and the exceptions of builtins. Saturation runs
within the prover's budget. Conditions decided by the path (\emph{this \code{if} is always
true here}) and redundant guards come out of the same saturation, extended per context.
Statement-level rules, such as a collecting loop that should be a comprehension or a handler
that swallows its exception, stay walks of the syntax.

Lint is opt-in (\code{check --lint}). Warnings that every run prints are noise, which agents
either chase or learn to ignore, and both are worse than not printing them. \code{--apply}
carries a lint out as a refactoring.

%=====================================================================
\section{Implementation}\label{sec:impl}
%=====================================================================

\subsection{Phases as packages}

The implementation is a package per phase, each importing only those before it, so the
imports enforce the layers (\cref{tab:size}). The driver \file{ulang/ulang.u} is a program
that imports them all and offers the commands \code{run}, \code{test}, \code{check},
\code{prove}, \code{compile}, and \code{format}. The runtime is about 1\,600 lines of C over
GMP and the Boehm collector.

\begin{table}[t]
\centering\small
\begin{tabular}{@{}lrl@{}}
\toprule
\textbf{Package} & \textbf{Lines} & \textbf{Contents} \\
\midrule
\file{ulang/syntax}  & 1\,247 & syntax, builtins, notation, errors shared by all phases \\
\file{ulang/format}  & 1\,103 & the canonical printer \\
\file{ulang/parse}   & 3\,078 & lexer, parser with recovery, loading of packages \\
\file{ulang/check}   & 5\,341 & signatures, inference, overloading, exhaustiveness, raises, messages \\
\file{ulang/run}     & 1\,560 & the interpreter \\
\file{ulang/compile} & 2\,057 & a typed syntax of C and the compiler \\
\file{ulang/prove}   & 2\,527 & terms, sequents, simplification, induction \\
\file{ulang/ulang.u} &   914 & the driver \\
\file{ulang/prelude} & 1\,522 & the library: lists, options, strings, sets, maps, arrays, buffers, os \\
\bottomrule
\end{tabular}
\caption{The self-hosting implementation (November 2026). The repository holds about 2\,070
named tests, among them 217 error cases, and 130 lemmas.}
\label{tab:size}
\end{table}

\subsection{Bootstrap}

The seed \file{bootstrap/ulang.c} is the C of the last compiler that compiled itself to the
same C. \code{make} builds the seed and translates the current sources with it, in about 15
seconds from nothing. \code{make selfcompile} replaces the seed only if the new compiler
reaches the fixed point. A rule that the seed's checker would reject in the new sources
(such as the exact comparison of raised types) takes two seeds. The way through is a copy of
the sources without the new annotations. That procedure is recorded, because it is exactly
the kind of step an agent would otherwise rediscover by failing.

\subsection{Backends}

The interpreter and the compiler share the checked program and the texts of runtime errors.
The compiler builds C as a small typed syntax with a printer that adds parentheses. Each
datatype gets one type descriptor with its order, equality, and show. Outs and refs are
copied in and out, \code{try} saves the variables its body assigns before \code{setjmp}, and
comprehensions over lists and arrays are fused into loops. Every value is a uniform
\code{u_value}, with no specialisation by type. That costs some speed, and it is what makes
separate compilation simple. The plan\planned\ compiles each package once into a static archive,
the library included, with mangled names that are stable across translation units and a stamp
symbol per archive, so that a stale archive fails to link rather than to run. A second
backend\planned\ writes x86-64 assembly from the same typed C, so that programs need only an assembler:
on the compiler itself, \code{cc -O2} takes 11.2\,s where \code{as} takes 0.2\,s, for a
measured slowdown of about 1.35$\times$ in a simple backend that keeps every variable in memory.

\subsection{Speed of the loop}

What matters for an agent is how long one round of feedback takes. Checking the whole
compiler takes about 0.9\,s, and the test suite of the project about 9\,s interpreted, with the
error cases compiled. The interpreter is about 300 times slower than compiled code. A faster
interpreter was rejected; the quick loop on large inputs is compiled tests instead, with one
tool command that builds and caches the binary of the tests\planned.

\subsection{Documentation that is checked}

The documentation is what an agent learns the language from, so it is tested like code.
\file{checks/docs.u} verifies that every reference in \file{CLAUDE.md}, \file{README.md}, and
\file{TODO.md}, written as a list of names in backticks followed by ``of'' and a file, names
declarations of that file. A rename that forgets the documentation fails the build. The
planned reference \file{LANGUAGE.md}\planned\ separates what a writer of programs needs (the
constructs by example, the semantics of effects, the library, and the differences from other
languages) from the implementation notes, which are loaded into every session today and
cost context.

%=====================================================================
\section{Status}\label{sec:status}
%=====================================================================

\begin{table}[h]
\centering\small
\begin{tabularx}{\textwidth}{@{}lX@{}}
\toprule
\textbf{Implemented} & core language, inference, overloading, hooks, forms and comprehensions,
pattern aliases, modules and packages, checked and inferred \code{raises}, atomic effects,
canonical format enforced by \code{check}, docs as syntax, positions on every part, all type
errors of a package in one run, parser resumes at the next command, interpreter, compiler to C,
compiled test suite, lemmas with the default proof, editor extension \\
\midrule
\textbf{Planned, with a plan} &
JSON diagnostics, exit statuses, recovery inside commands, hints (\file{diagnostics.md});
examples and enumerated lemma tests (\file{examples.md});
e-graph core, budgets, decision procedures, proof language (\file{prover.md}, \file{prove.md});
lint by extraction (\file{lint.md});
archives per package (\file{archives.md}); assembly backend (\file{assembly.md}) \\
\midrule
\textbf{Open, in \file{TODO.md}} &
\code{resolve}/\code{refs}/\code{sig}; refactorings; tests with formals; \file{LANGUAGE.md};
exceptions as values; termination check; unused imports as errors; stopping folds and bounded
quantifiers; qualified names; \code{os} as a package; a single build tool; positions in
runtime errors; contracts and \code{correct}; evidence per obligation; \code{uses os};
replay of \code{os} results \\
\bottomrule
\end{tabularx}
\caption{Status of the parts described in this paper.}
\end{table}

%=====================================================================
\section{Reconsidering the Design}\label{sec:discussion}
%=====================================================================

With the plans assumed done, this section asks where the design is still in tension with its
own premises, and what we would change or test.

\subsection{What the design gets right}

\paragraph{The function/procedure split is the cheapest effect system that pays.}
One bit per declaration, visible in its keyword, gives most of the value of effect rows: a
function cannot read a file, raise, or loop forever (once termination is checked). Checked
exceptions with inference for private procedures put the annotation burden only at
boundaries, where readers need it. \code{uses os} on exported procedures\planned\ would add
the remaining capability at the same price.

\paragraph{Atomic effects remove a whole class of review.}
In most languages, reviewing a handler means asking which assignments in the \code{try}
happened before the exception. In ulang the answer is always ``none''. The plain reading of
\cref{fig:bank} is the correct one, and an agent reasoning locally reaches it.

\paragraph{Uniqueness without preference is the right overloading for agents.}
Preferences (``the most specific'', ``the nearest'') are rules that a reader must simulate.
ulang's rule needs no simulation: a use compiles only if exactly one candidate fits, and
otherwise the error lists the candidates. The planned \code{resolve} makes the remaining
invisibility a single lookup.

\paragraph{Canonical form and docs as syntax compose.}
Because the layout is canonical and docs are nodes, every tool that changes the tree
(refactorings, \code{--apply} of lints, examples printed back) can print the whole file
without losing anything. This is the precondition for typed edits, and few languages meet it.

\subsection{Tensions that remain}

\paragraph{Novel syntax against zero training data.}
ulang deliberately breaks habits (\code{x := e}, \code{[1, 2]}, \code{+} on strings,
\code{return}). Each break is a learning cost for an agent with no ulang in its priors. The
bet is that explicit, uniform syntax plus hints (P11) costs less than familiar but ambiguous
syntax. We have anecdotes: \file{unclear.md} records about thirty trip-ups, mostly in the
implementation rather than the language, and the habits of \cref{tab:hints} are those
observed. We do not have a measurement. \emph{Proposed experiment:} a fixed set of tasks, given
to agents with only \file{LANGUAGE.md}, measuring the rounds to a passing program with and
without hints, and against a variant with familiar spellings for the same semantics. If the
familiar variant wins clearly, the forms and keywords are worth revisiting. The semantics,
which are where the value lies, would stay.

\paragraph{Locality has three exceptions.}
User-defined notations make parsing depend on imports. Cross-module overloading makes
resolution depend on imports. Hooks make \code{==} on a type depend on a function in its
module. Each was kept for a reason: notations for the prelude's operators, overloading for
shared container names, hooks for abstract types. The hooks are the most worrying, since a
declared \code{compare} must agree with structural equality, and nothing checks that today.
\emph{Recommendation:} generate the agreement as a lemma, tested on enumerated values and
proven where possible, and make the hooks a line of \code{sig} output, so that they are seen
where the type is.

\paragraph{Argument order is the residual grounding hazard.}
Named fields protect constructors, and types protect most calls. They do not protect
\code{slice(i, j, s)} or \code{split(sep, s)}. Named arguments at calls would protect them,
but they add a second way to write every call, against P2. The design's answer is examples
in signatures, plus a hint that tries permutations of the arguments. We think that is right,
and suggest making it a style rule with a check: an exported function with two formals of
the same type must have an example.

\paragraph{Exceptions cannot become values.}
Code that collects errors (the driver, the test helpers) catches each kind into an out and
lists every phase in long \code{raises}. This is the most visible awkwardness in the sources.
A \code{try} with a result, or a statement that gives an \code{Option} of the exception of a
call, together with the errors of all phases as one datatype, would fix it without weakening
checked exceptions. We would do this before the contracts, since the prover's obligations for
\code{raise} paths become simpler with it.

\paragraph{Partial field accessors.}
\code{e.f} raises \code{MissingField} for a constructor without the field: a partial function
in a language that otherwise makes failure visible. Marking partial access or rejecting it
statically were both considered and dropped. With \code{requires} on functions and the checker
knowing which constructors have a field, the accessor could carry an implicit
\code{requires e is C1 or ...}, turning the defect into an obligation, consistent with the
treatment of \code{/}.

\paragraph{The library is ambient.}
Every module sees the whole library, including \code{os}, without imports. That keeps short
programs short, but it makes every library name a candidate everywhere and hides I/O.
Moving \code{os} into its own package, which must be imported, is planned and costs a line
per program that does I/O. We consider that line a feature, as a declaration of capability.

\paragraph{Verbosity has a reading cost after all.}
P3 says keystrokes are cheap. But context windows are finite, and \file{CLAUDE.md}, which
describes the language, its implementation, and the workflow in sentences that carry a
reference in every clause, now has about 9\,000 words, loaded into every session. The fix is
not a terser language. It is separating what writers of programs need (\file{LANGUAGE.md})
from what implementers need, and measuring the paragraph cost of each feature, as P1--P11
suggest.

\paragraph{Budgets make proofs deterministic, not robust.}
A proof that closes at one budget may stay open at a smaller one, and a stronger simplifier
can only close more goals, never fewer. That holds only if extraction and saturation are
monotone in the facts, which an e-graph gives, but a budget of nodes cuts across. Reporting
exhaustion is necessary. A test that runs the corpus of lemmas at half the budget and lists
the ones that open would show how much headroom proofs have.

\paragraph{Functions without parameters are computed once.}
They are effectively global constants, initialised lazily. That is safe because functions are
pure, but a non-terminating one (until termination is checked) deadlocks its first use.
The termination check closes this.

\subsection{What we would add, beyond the plans}

\begin{itemize}
\item \textbf{A rejected-habits test suite as a specification of P11}: one case per habit,
  with the hint expected, so that the documentation and the checker cannot drift apart.
\item \textbf{A paragraph budget}: each feature's explanation in \file{LANGUAGE.md} is
  measured, and a new feature states its paragraph before it is implemented, as the cost
  P1--P11 accept.
\item \textbf{Evidence in the output}: \code{test} and \code{prove} report, per declaration,
  the highest rung of \cref{tab:ladder} it reaches, so that an agent (or a reviewer of an
  agent's change) sees which claims are proven, which are only tested, and which are open.
\item \textbf{Replay of runs}: \code{run} records the results of \code{os} procedures and
  replays them, so that a failure an agent observed can be reproduced by a human exactly.
  This follows from P9 and costs little, since all I/O goes through one module.
\end{itemize}

%=====================================================================
\section{Related Work}\label{sec:related}
%=====================================================================

\paragraph{Modes and contracts.}
The \code{in}/\code{out}/\code{ref} modes come from Ada and SPARK~\cite{barnes2003spark},
whose discipline of data flow ulang shares. Contracts follow Eiffel~\cite{meyer1997oosc} and
the auto-active verifiers Dafny~\cite{leino2010dafny} and Why3~\cite{filliatre2013why3}.
ulang differs in keeping proofs and loop specifications out of the code, in using no external
solver, and in putting enumerated tests below proofs on the same statements.

\paragraph{Effects.}
Koka~\cite{leijen2014koka} tracks effects as rows. ulang's function/procedure split with
checked, inferred exception sets is a much coarser point in that space, chosen for its
paragraph cost. Java's checked exceptions~\cite{gosling1996} are the cautionary precedent.
ulang avoids their main burden by inferring the raises of private procedures and by having no
subtyping among exceptions, so the sets compare exactly.

\paragraph{Uniformity and canonical form.}
Go~\cite{donovan2015go} made a canonical formatter part of the culture; ulang makes it part
of the checker. Unison~\cite{unison} identifies definitions by the hash of their syntax. ulang
uses the e-graph of the prover for the weaker but semantic goal of \emph{detecting}
duplicates and simpler equivalents, not of identifying them.

\paragraph{Proving.}
The simplifier follows the tradition of Isabelle's \code{simp}~\cite{nipkow2002isabelle} and
of cooperating decision procedures~\cite{nelson1979,shostak1984,conchon2008ccx}. The e-graph
plan follows egg~\cite{willsey2021egg}. Higher-order patterns follow Miller~\cite{miller1991},
and the structured proof language is in the spirit of Isar and of Lean~4's tactic
blocks~\cite{demoura2021lean4}, reduced to a handful of keywords. Enumerated testing follows
SmallCheck~\cite{runciman2008smallcheck} rather than QuickCheck~\cite{claessen2000quickcheck},
for determinism.

\paragraph{Languages for agents.}
We are not aware of a language designed from the premise that its authors learn it in
context and its main cost is verification. Most work on agents and code adapts the
environment (tools, prompts, retrieval) to existing languages. ulang asks the converse
question: what the language should be when the environment is an agent.

%=====================================================================
\section{Conclusion}
%=====================================================================

ulang shows that a language designed around cost to verify and in-context learnability is
practical: small enough to explain in a few pages, strict enough that the checker teaches
it, and expressive enough to implement its own compiler and prover. Its central choices are
pure functions and effectful procedures, modes at every call, atomic effects, checked and
inferred exceptions, overloading without preferences, one way to write each construct, a
canonical layout, and docs as syntax. Each one removes a question that a reviewer of an
agent's code would otherwise have to ask. The plans now in progress make the checker's feedback a
protocol with hints and positions on every part, put examples into signatures, test lemmas
before proving them, and derive lint from proven equalities. The largest open question is
empirical: whether the departures from familiar syntax pay for themselves with agents that
have never seen the language. The experiment of \cref{sec:discussion} would answer it, and
the design is meant to change according to the result.

\paragraph{Acknowledgement.}
The implementation of ulang has been developed largely with coding agents, whose trip-ups,
recorded in \file{unclear.md}, informed many of the decisions described here.

\paragraph{Declaration of AI use.}
This paper was drafted by an AI assistant (Claude Opus 5.5, in Claude Code), working from
the repository of ulang: its sources, \file{CLAUDE.md}, \file{README.md}, \file{DESIGN.md},
\file{TODO.md}, and the plans in progress. The assistant also ran the code examples and
error messages quoted here against the current implementation. The design itself, its
principles, and the plans the paper describes are the author's, developed together with
coding agents. The author is responsible for the content, which remains to be reviewed and
revised.

\bibliographystyle{plainnat}
\bibliography{refs}

\end{document}
